% Part: lambda-calculus
% Chapter: lambda-definability
% Section: pairs

\documentclass[../../../include/open-logic-section]{subfiles}

\begin{document}

\olfileid[fr]{lam}{rep}{pai}
\olsection{Couples et prédécesseur}

\begin{defn}
Le couple de $M$ et $N$, noté
$\tuple{M,N}$, est défini par
\[
  \tuple{M,N} \ident \lambd[f][fMN],\quad f\notin\FV{M}\cup\FV{N}.
\]
\end{defn}

Intuitivement, ce terme reçoit une fonction
et l'applique aux deux composantes du
couple. Le constructeur suivant reçoit
deux termes et renvoie le couple qu'ils
forment :
\[
  \fn{Pair} \ident \lambd[mn][\lambd[f][fmn]].
\]

Pour retrouver les composantes d'un couple,
utilisons deux fonctions de projection :
\begin{align*}
  \fn{Fst} & \ident
    \lambd[p][p(\lambd[mn][m])]\\
  \fn{Snd} & \ident
    \lambd[p][p(\lambd[mn][n])].
\end{align*}

\begin{prob}
Expliquer pourquoi $\fn{Fst}$ et
$\fn{Snd}$ renvoient respectivement
la première et la seconde composante.
\end{prob}

Les couples permettent de $\lambda$-définir
la fonction prédécesseur. Pour rendre
lisible le terme de la source, nommons~$J$
la fonction qui met à jour un couple :
\begin{align*}
  J & \ident
    \lambd[p][\tuple{\fn{Snd}\,p,
      \fn{Succ}(\fn{Snd}\,p)}],\\
  \fn{Pred} & \ident
    \lambd[n][\fn{Fst}(n J
      \tuple{\num{0},\num{0}})].
\end{align*}
Cette écriture abrège exactement le terme
de la source, où la définition de~$J$
est insérée à la place de son nom.

Rappelons que $\num{n}\,f x$ se réduit
à $f^n(x)$. Ici, $f=J$ reçoit un
couple~$p$ et renvoie celui qui contient
en première composante la seconde
composante de~$p$, puis en seconde
composante le successeur de celle-ci.
Le point de départ est
$\tuple{\num{0},\num{0}}$.
Après zéro itération, le résultat est
ce couple initial. Après $n>0$
itérations, c'est
$\tuple{\num{n-1},\num{n}}$.
La projection $\fn{Fst}$ en extrait
le prédécesseur, avec
$\fn{Pred}(\num{0})\red\num{0}$.

La soustraction s'obtient en appliquant
$\fn{Pred}$ $b$ fois à~$a$ :
\[
  \fn{Sub} \ident
  \lambd[ab][b \fn{Pred}\, a].
\]
Sur les naturels, cette soustraction
est tronquée à zéro lorsque $b>a$.

L'invariant de l'itération peut s'écrire
\[
J^n\tuple{\num{0},\num{0}}
\red \tuple{\num{\max(n-1,0)},\num{n}}.
\]
Le couple initial vérifie la formule pour $n=0$.
La mise à jour $J$ place ensuite le deuxième numéral
dans la première composante et son successeur dans la seconde.

\end{document}
